%
% File acl2020.tex
%
%% Based on the style files for ACL 2020, which were
%% Based on the style files for ACL 2018, NAACL 2018/19, which were
%% Based on the style files for ACL-2015, with some improvements
%%  taken from the NAACL-2016 style
%% Based on the style files for ACL-2014, which were, in turn,
%% based on ACL-2013, ACL-2012, ACL-2011, ACL-2010, ACL-IJCNLP-2009,
%% EACL-2009, IJCNLP-2008...
%% Based on the style files for EACL 2006 by 
%%e.agirre@ehu.es or Sergi.Balari@uab.es
%% and that of ACL 08 by Joakim Nivre and Noah Smith

\documentclass[11pt,a4paper]{article}
\usepackage[hyperref]{acl2020}
\usepackage{times}
\usepackage{latexsym}
\usepackage{graphicx}
\usepackage{float}
\usepackage[T1]{fontenc}
\renewcommand{\UrlFont}{\ttfamily\small}

% This is not strictly necessary, and may be commented out,
% but it will improve the layout of the manuscript,
% and will typically save some space.
\usepackage{microtype}

\aclfinalcopy % Uncomment this line for the final submission
%\def\aclpaperid{***} %  Enter the acl Paper ID here

%\setlength\titlebox{5cm}
% You can expand the titlebox if you need extra space
% to show all the authors. Please do not make the titlebox
% smaller than 5cm (the original size); we will check this
% in the camera-ready version and ask you to change it back.

\newcommand\BibTeX{B\textsc{ib}\TeX}

\title{Exploring Norma for Historical Text Normalization}

\author{Tova Erbén  \\
  LT2314 Language Technology Resources \\
  Gothenburg University \\
  \texttt{guserbto@student.gu.se} }

\date{}

\begin{document}
\maketitle
\begin{abstract}
This work aims to explore the Norma tool for historical text normalization using substitution lists. The material used in this endeavour is the MAÞiR resource from Språkbanken, containing annotated excerpts from five different texts in Old Swedish (ca 1225-1526).
\end{abstract}

\section{Introduction}

Historical texts that were written before any spelling standardization had taken place in their language of medium tend to contain a lot of spelling variation. This becomes problematic when you want to parse them for further processing or research them in a corpus \citet{abc}. There exists several different methods for normalizing such text, i.e. mapping them to a standard form. 

\citet{bollmann:19} gives an overview and comparison of the following methods: 1) substitution lists, 2) rule-based methods, 3) distance-based methods, 4) statistical machine translation, and 5) neural machine translation. 

Substitution lists (also known as lookup-based normalization) are the simplest out of these methods and consists of putting together a list of different variants for different words. Variants that don't make it to the list cannot be handled. The normalization tools VARD and Norma contain such substitution list components.

Rule-based methods instead try to pick up on regularities in the variations, such as the letter \emph{v} often being written like \emph{u}. \citet{bollmann:12} describe the rule-based re-write rules as operating on one or more characters and only taking the immediate character context into consideration. If more than one rule is applicable, then you choose the one with the highest frequency in the training data.   

The distance-based methods look at the number of edits needed to go from a variant spelling to a standard form, often using the Levenshtein distance. \citet{adesam:12} combine the distance- and rule-based approach by using the Levenshtein distance to derive rewrite rules for Old Swedish, which perform better when compared to manually compiled rewrite rules from dictionaries.

The statistical machine translation-based methods look to optimize the probability for a spelling variant \emph{v} being of standard form \emph{w} by looking, for instance, at character-based statistics (cSMT). \citet{pettersson:13} manage to increase the normalization accuracy for Swedish from 64.6 to 92.3 \% by applying cSMT to historical texts from the time period 1527–1812.    

A recent data-driven approach for spelling normalization is neural machine translation (NMT), commonly using encoder-decoder models with LSTM units \citep{bollmann:19}. However, when \citet{bollmann:19} compares some character-based NMT models to the previously mentioned methods, he notes that these models need more data in order to perform well. Instead, he recommends using the Norma tool in the ``combined'' setting (more about this later) for languages with little training data, and the character-based statistical machine translation tool cSMTiser otherwise.     
\section{Material and Software}

\subsection{MAÞiR}
The material that I will be using in this study comes from the MAÞiR resource at Språkbanken, consisting of annotated excerpts from five different Old Swedish texts:

\begin{enumerate}
   \item \emph{Abota}
   \begin{itemize}
     \item ca 1448
     \item 540 tokens
   \end{itemize}
   \item \emph{Avgl} 
     \begin{itemize}
     \item Elder Westrogothic law
     \item ca 1220/1280
     \item 1228 tokens
   \end{itemize}
    \item \emph{Moses-b} 
     \begin{itemize}
     \item Paraphrases of the five books of Moses 
     \item ca 1330/1526
     \item 7049 tokens
   \end{itemize}
    \item \emph{Ogl-a} 
     \begin{itemize}
     \item Ostrogothic law
     \item ca 1290/1350
     \item 22090 tokens
   \end{itemize}
   \item \emph{Tungulus} 
     \begin{itemize}
     \item 1457
     \item 2814 tokens
   \end{itemize}
\end{enumerate}

These texts have been annotated word-by-word in XML format with information about syntactic relations, part of speech tags, morphology, lemma and head. This part of the resource is called MaÞiR Trees. The other one, MaÞiR Words, contains 10 TSV ('tab-separated values') files with lemmata information extracted from a reference dictionary for Medieval Swedish by \citet{soderwall:1891-1918} (Sdw). 

Half of the files have four columns with information about lemma (which may contain spaces, and in the case of the verb file, a '+' sign separating the verb from a separable participle, e.g. EX), pseudolemma (containing only the verbal head in the case of verbs with separable participles), coarse-grained POS tag in accordance with Sdw, and finally a fine-grained POS tag. The other five files contain an additional column with parenthetical form, which show inflectional and spelling-related variation, also taken from Sdw.    

\begin{figure}[H]
    \centering
    \includegraphics[width=7cm]{images/tree.jpg}
    \caption{my tree}
    \label{fig:mytree}
\end{figure}

\subsection{Norma} 

Norma\footnote{https://github.com/comphist/norma} \citep{bollmann:12} is an automatic normalization tool that contains functionality for rule-based, distance-based and lookup up-based normalization, as well as the option of running all these three in combination. When running the program, you have to provide it with a configuration file that specifies which normalizer(s) to use, and in which order. 

When the program reads the first historical input word to be normalized (given in an input text file with one word per line), it proceeds to find a normalized candidate through the first normalizer specified in the configuration file. If the candidate word is above a certain confidence score threshold (between 0 and 1) defined by the user, it moves on to the next step of validating the candidate against the correct word form and starting the training function which takes the historical spelling and the standardized word form as input and adjusts the parameters of the normalizer. If the first normalizer does not yield a high enough confidence score, it moves on to the next normalizer. If none of the specified normalizers give a high enough score, the word remains unchanged.         

Norma has three modes: interactive mode, training mode and batch mode. The interactive mode is supposed to allow the user to confirm or correct the normalization candidate provided by Norma themself, but is currently unavailable in the latest release of Norma due to technical issues. The batch mode is the default mode which takes the configuration file and text file of words to normalize as input and outputs a list of normalized words, one token per line together with a confidence score. The training mode allows you to provide manually normalized data in the form of a list with two tokens per line: historical form and the normalized form. 

According to the user guide\footnote{https://github.com/comphist/norma/blob/master/doc/UserGuide.md}, "most normalizers" require a target lexicon that defines valid output strings in order to avoid generating non-existing words. This, naturally, does not seem to be the case for the Mapper normalizer, as it only normalizes words to forms that exist in the substitution list. 

\section{Method}
For preprocessing the annotated MAÞiR files I wrote a Python script that takes the path to the XML files as argument and then loops them through in order to extract tokens and their lemma. These are saved in two files: 

\begin{enumerate}
   \item \texttt{mathir\_tokens.txt}, which lists all tokens (including duplicates) from all files in lower case and one token per line:
   \begin{table}[H]
    \centering
    \begin{tabular}{c}
       hær \\
       sigx \\
       aff \\
       abotum \\
       allum \\
       \ldots
    \end{tabular}
   \end{table}
    \item \texttt{mathir\_train.txt}, which lists all tokens and their corresponding lemma, one tab-separated token-lemma pair per line:
    
    \begin{table}[H]
    \centering
    \begin{tabular}{c|c}
        hær & här \\
        sigx & sighia \\
        aff & af \\
        abotum & abbote \\ 
        allum &	alder \\ 
        \ldots & \ldots
    \end{tabular}
    \end{table}
\end{enumerate}

This latter file is only used for calculating the baseline accuracy, i.e. the number of tokens that already match their corresponding lemma.

As basis for a subsitution list, I instead use five of the files in MAÞiR Words which contain parenthetical form or variant spelling of the dictionary lemmas. Therefore I wrote another script that takes the MAÞiR Words directory as input and loops through these files, saving the parenthetical form and corresponding lemma in another text file called \texttt{sdw\_train.txt}: 

    \begin{table}[H]
    \centering
    \begin{tabular}{c|c}
        \ldots & \ldots \\
        lach & lagh \\
        lagha & lagh \\
        laghä & lagh \\
        lak & lagh \\ 
        laugh &	lagh \\ 
        \ldots & \ldots
    \end{tabular}
    \end{table}

When running this file in train mode in Norma, a parameter file called \texttt{sdw\_train.Mapper.mapfile} is automatically generated. In those cases where the training data is not coming from dictionary entries (like here), but instead generated from real data like the \texttt{mathir\_train.txt} file, this file will contain frequencies of each variant-lemma pair, which will then help Norma determine which lemma to choose when doing the lookup.

For normalizing the tokens, I simply ran the \texttt{mathir\_tokens.txt} file in batch mode and saved the output to a text file by adding \texttt{>norma\_output.txt} to the Norma command.     

For evaluating the output file, I wrote a script that takes the token file and output file as arguments and counts the number of matches between these two documents. The accuracy is calculated by simply dividing the total number of matches with the total number of tokens.

\section{Results}

    \begin{table}[H]
    \centering
    \begin{tabular}{c|c}
        \hline
        \textbf{Normalizer} & \textbf{Accuracy} \\
         \hline
        Baseline (unchanged) & 27.30 \% \\
        \hline
        Norma (lookup) & 34.14 \%  \\
        \hline
    \end{tabular}
    \caption{\label{tab:table-results}Accuracy for baseline (leaving all tokens unchanged) and look-up based normalization in Norma.}
    \end{table}

The lookup-based normalization in Norma yielded a lemmatization accuracy of 34.14 \%, which can be compared to the baseline of 27.30 \%. I would have wanted to try out the weighted Levenshtein distance-based normalization in Norma as well, but unfortunately the command for generating the necessary target dictionary files was not working in Docker. See the figure below for examples of variant-lemma pairs that were missing in Sdw.

    \begin{table}[H]
    \centering
    \begin{tabular}{c|c}
        \textbf{Matches} & \textbf{Failed Matches} \\
        \textbf{input} \quad \textbf{output} & \textbf{input} \quad \textbf{output} \quad \textbf{correct} \\
         \ldots & \ldots \\
        oc  \quad ok & wt \quad wt \quad ut \\
        uili \quad vilia & ær \quad ær \quad vara \\
        scal \quad skula & thet \quad tea \quad þän \\
        med \quad mäþ & æter \quad æter \quad äta \\
        þessi \quad þänne & þa \quad þiggia \quad þa \\
        prester \quad präster & bætri \quad bætri \quad bätre \\
        bönder \quad bonde & sigher \quad sigha \quad sighia \\
        ligger \quad liggia & wigyæ \quad wigyæ \quad vighia\\
        kyrkiu  \quad kirkia & guðfæþur \quad guðfæþur \quad guþfaþir\\
        landbor \quad landboe &	for budhit\quad for budhit \quad forbiuþa \\
          himmerike \quad himirike & wtan widh \quad wtan widh \quad utan viþer \\
        \ldots & \ldots
    \end{tabular}
    \end{table}

\section{Discussion}



\section{Conclusion and further work}

\bibliography{anthology,acl2020}
\bibliographystyle{acl_natbib}

\end{document}
